\section{第~\ref{sec:livingmachine}~章　用活神经元搭成的机器}

电极阵列上培养物的行为是在带记录和刺激的直接实验中测得的；群放电机制依据突触耗竭模型和微培养物实验，而关于培养皿中多巴胺的论断依据的是个别实验，其中一部分连作者本人也只视为演示。

{\footnotesize
\setlength{\tabcolsep}{3pt}\renewcommand{\arraystretch}{1.15}
\begin{longtable}{>{\raggedright}p{0.5cm}>{\raggedright}p{4.0cm}>{\raggedright}p{5.6cm}>{\raggedright}p{2.3cm}>{\raggedright\arraybackslash}p{1.7cm}}
\toprule
序号 & 论断 & 实验及测得的结果 & 来源 & 状态 \\
\midrule\endfirsthead
\toprule
序号 & 论断 & 实验及测得的结果 & 来源 & 状态 \\
\midrule\endhead
1 & 芯片上的培养物：主要是\nt{GLUT}，较小比例的\nt{GABA}，比例取决于标本；海马培养物中约 $6\%$ & 电极阵列上由几千个神经元组成的网络是标准实验（第~3~行）。海马培养物中\nt{GABA}神经元的比例用GABA染色计数：约占神经元的 $6\%$。对皮层培养物，我们不给出经过核实的数字 & \cite{Benson1994} & 已测量 \\
2 & 类器官：由人类干细胞培养、大小约一毫米的三维神经组织 & 用干细胞培养出含有多个脑区的三维类器官，其中包括带有前体细胞层和成熟神经元的皮层 & \cite{Lancaster2013} & 已测量 \\
3 & 无输入时培养物以群放电形式爆发，间隔从一秒到几分钟，因培养物而异 & 对密度为 $3\,000$ 到 $50\,000$ 个神经元的 $58$ 个皮层培养物记录了五周（$963$ 段半小时记录）：两周后几乎所有培养物都以全网络群放电为主；群放电间隔为 $1$ 到 $300$~s，并强烈依赖于标本 & \cite{Wagenaar2006} & 已测量 \\
4 & 群放电因递质耗竭而结束，间隔 $T_{\text{burst}}\approx\tau_{\text{rec}}\ln\frac{1}{1-x^*}$~\eqref{eq:burstinterval}；$2$--$4$~s是 $\tau_{\text{rec}}=0.8$~s时的模型估计，实测的恢复更慢 & 公式在本章推导。模型：带耗竭性突触的网络会自发产生全网络群放电。在 $4$--$30$ 个神经元的微培养物上的实验：群放电时长取决于距上次群放电的时间，递质囊泡耗竭会取消群放电；作者的结论是群放电受递质储备限制。那里储备恢复需要几十秒，代入同一公式得到几十秒乃至几分钟的间隔 & 本章推导；\cite{Tsodyks2000,CohenSegal2011,Tsodyks1997} & 理论 \\
5 & “会闭嘴的老师”：出现所需响应时关掉刺激，响应被固定下来 & 以 $0.3$--$1$~Hz刺激培养物，直到刺激后 $50\pm10$~ms出现选定的响应，然后关掉刺激 $5$~min；几个循环后所需响应立刻出现；绘制了学习曲线 & \cite{Shahaf2001} & 已测量 \\
6 & 培养物打网球；一次会话内连续击球回合的显著增长只出现在完整反馈下 & 详见第~\ref{sec:dishbrain}~章及其补充：以安静代替刺激时，培养物到会话结束时也比无反馈时好，但效应较小；跨会话学习不稳定 & \cite{Kagan2022} & 已测量 \\
7 & 培养物学会预测自己的输入是假说；关于其“有感知能力”的高调说法遭到反驳 & 自由能原理是一个理论设想；没有直接实验能把它与更简单的解释区分开。三十位研究者在回应文章中指出，闭环中的行为变化既不表明智能，也不表明感知 & \cite{Friston2010,Balci2023} & 假说 \\
8 & 在挑选好的刺激模式下，培养物驾驶虚拟身体走向目标 & 程序根据当前结果自行挑选训练刺激模式；几十分钟内网络达到指定状态，训练 $2$~h后可塑性保持 $80$~min。同样的刺激若与行为无关地施加，则不起作用 & \cite{Bakkum2008} & 已测量 \\
9 & 储备池：只训练线性读出 $y=W_{\text{out}}\mathbf r$~\eqref{eq:reservoir} & 储备池计算理论：任意带衰减记忆的非线性系统加上一个训练好的线性输出 & \cite{Maass2002,JaegerHaas2004} & 理论 \\
10 & 作为储备池的类器官辨别说话人的准确率约 $78\%$（作者数据）；按误差学习的是读出，类器官的连接在无监督地变化 & 把几位说话人的语音以电极阵列上的电流模式输入类器官；训练好的读出辨别说话人的准确率约 $78\%$，这是作者报告的。作者还报告训练期间类器官的功能连接发生了变化，称之为类器官内部的无监督学习 & \cite{Cai2023} & 已测量 \\
11 & 一百万个神经元用于脉冲的功耗约 $10^{-4}$~W~\eqref{eq:dishpower} & 估计：由皮层能量预算~\eqref{eq:energy} 得出的每个脉冲 $10^{-10}$~J，乘以脉冲数；培养物功率没有直接测量 & 本章推导；\cite{Attwell2001} & 理论 \\
12 & 闪光释放多巴胺：$1$~mM笼锁多巴胺，$240$ 次重复，诱发脉冲数增加 & 在远程类器官平台上，当刺激引发更多脉冲时，用紫外光（$365$~nm，$0.8$~s）释放光敏笼中的多巴胺（$1$~mM）；$240$ 步（约 $5$~h）中脉冲数总体增加。作者写道，仅凭一次实验无法确立与多巴胺的关联；没有对照 & \cite{Jordan2024} & 假说 \\
13 & 活细胞分泌的多巴胺能长期、可逆地改变有机晶体管的权重 & 分泌多巴胺的细胞培养在有机神经形态元件上方；多巴胺在电极处氧化改变沟道电导；展示了权重的长期改变及其复原 & \cite{Keene2020} & 已测量 \\
14 & 存在含有工作\nt{DOPA}神经元的类器官；其多巴胺未被用作老师 & 在人类干细胞类器官中发现了电活动正常的成熟\nt{DOPA}神经元以及多巴胺的产生。文献中我们没有找到用这种多巴胺去教另一个网络的实验 & \cite{Jo2016} & 已测量 \\
15 & 类器官平衡杆子：程序给出的教学刺激优于随机刺激，不能巩固，经由\nt{GLUT}突触 & 小鼠皮层类器官与“小车倒立摆”任务构成闭环：对大多数类器官，强化学习选择的信号比随机信号或无信号效果更好；休息 $45$~min后改进不保留；阻断AMPA和NMDA受体则改进消失 & \cite{Robbins2026} & 已测量 \\
16 & 没有控制寄存器，试验台上的网络无法自己决定什么算成功（命题~\ref{prop:livingmachine}） & 在第~5--15~行的所有实验中，教学信号都由实验者或程序设定；没有培养物自己形成成功标准的实验——这是从“不存在”得出的推论，不是实验 & --- & 假说 \\
\bottomrule
\end{longtable}}
